% Generated draft: requires mathematical and source review.
\documentclass[11pt]{article}
\usepackage[margin=25mm]{geometry}
\usepackage{fontspec}
\setmainfont{Latin Modern Roman}
\usepackage{amsmath,amssymb,mathtools,unicode-math}
\setmathfont{Latin Modern Math}
\usepackage{enumitem,needspace,xurl,hyperref}
\setlength{\parindent}{0pt}
\setlength{\parskip}{5pt}
\setlength{\emergencystretch}{3em}
\newcommand{\N}{\mathbb{N}}
\newcommand{\Z}{\mathbb{Z}}
\newcommand{\Q}{\mathbb{Q}}
\newcommand{\R}{\mathbb{R}}
\newcommand{\C}{\mathbb{C}}
\newcommand{\F}{\mathbb{F}}
\newcommand{\A}{\mathbb{A}}
\newcommand{\PP}{\mathbb{P}}
\newcommand{\E}{\mathbb{E}}
\newcommand{\eps}{\varepsilon}
\newcommand{\dd}{\,\mathrm{d}}
\begin{document}
\section*{UnsolvedMath ID: 30006505}
\textbf{Source problem code:} catalogue entry
\subsection{Definitions and mathematical statement}
Let \(X\) denote a space whose microscopic and macroscopic geometries differ, as described qualitatively in the supplied source record.

\textbf{Scale-detection problem.} Develop an intrinsic procedure that detects the appropriate scale at which to measure volume, perimeter, and isoperimetric behavior on \(X\).

\textbf{Unresolved formulation.} The supplied record does not specify a class of admissible spaces, the structures given on \(X\), definitions of volume or perimeter, or a mathematical criterion for distinguishing microscopic from macroscopic geometry. It also does not define an intrinsic procedure, its required output, or what makes a measurement scale correct. In particular, it is unspecified whether a single scale should serve all three measurements, whether scales may depend on location or on the set being measured, and which isoperimetric inequality or scale-dependent behavior is sought. Consequently, this is a qualitative research question; the supplied evidence does not determine a fully quantified mathematical assertion.

The extracted original text mentions fractal spaces and includes a figure caption describing an approximation of a Sierpinski gasket cable system. This supplies context, but does not establish that such cable systems are the intended exclusive domain or provide additional hypotheses.
\subsection{Short English statement}
Find an intrinsic way to identify suitable measurement scales in spaces with different small-scale and large-scale geometries, particularly in the fractal setting mentioned in the source. The admissible spaces, measurement definitions, and criterion for a correct scale remain unspecified, so a precise theorem statement requires further source clarification.
\subsection{Sources}
\begin{itemize}
\item Supplied local registry record: permanent UnsolvedMath ID \(30006505\), title \emph{Detecting Geometric Scales in Fractal Measurements}, source problem code \texttt{catalogue entry}, in the Oberwolfach Reports Open Problems collection. The registry records proposed year \(2025\); its proposer field is empty.
\item Supplied report citation: \emph{Functional Inequalities: Geometric Calculus meets Stochastic Analysis} (2026), Oberwolfach Reports, DOI \texttt{10.4171/owr/2025/54}, pp.\ 2917--2994. \url{https://doi.org/10.4171/owr/2025/54}. The supplied original extraction associates the question with Patricia Alonso Ruiz, joint work with Fabrice Baudoin, under the title \emph{An isoperimetric inequality for (fractal) spaces with different micro- and macrostructure}. This attribution does not establish a separately identified problem proposer.
\item Supplied literature source: Alonso Ruiz--Baudoin, \emph{Besov classes on Dirichlet spaces}. \url{https://arxiv.org/abs/1903.10078}.
\item Literature assessment dated exactly \texttt{2026-08-22}: the supplied curation labels the problem open and reports that existing fractal-space isoperimetric results require specified structures and scales, and that no general theorem answering the extracted scale-detection question was found. The supplied checked-source list consists of Oberwolfach Report 2025/54 and the Alonso Ruiz--Baudoin reference above. This is supplied dated literature triage, not independently verified here; time-sensitive status and cited claims require independent verification.
\item The registry labels the statement \texttt{corrected\_verified} and says that only clear formula delimiters, notation, grammar, and typography were repaired while mathematical content and scope were preserved. These are supplied source-processing claims, not a mathematical verification of this draft.
\end{itemize}
\end{document}
